% Fixed-population analysis. Numerical inputs: generated/ch3_transport_report.json.
\chapter{Horizontal transport and displacement scaling}
\label{ch:fixed-transport}
% Generated by scripts/tesi/ch03_fixed_transport/write_ch3_text.py
% Chapter prose is hand-written; only these values are generated.
\newcommand{\ChThreeClosedAltShiftMax}{0.051}
\newcommand{\ChThreeClosedAltShiftMin}{0.021}
\newcommand{\ChThreeCohortBiasMaxPct}{2.0}
\newcommand{\ChThreeCohortBiasMinPct}{1.8}
\newcommand{\ChThreeCohortShiftMax}{0.0177}
\newcommand{\ChThreeCohortShiftMin}{0.0164}
\newcommand{\ChThreeEligibleClosed}{87,807}
\newcommand{\ChThreeEligibleClosedPct}{56.6}
\newcommand{\ChThreeEligibleOpen}{67,003}
\newcommand{\ChThreeEligibleOpenPct}{43.2}
\newcommand{\ChThreeEligibleUnknown}{275}
\newcommand{\ChThreeGeneralFitLagMax}{30000}
\newcommand{\ChThreeGeneralFitLagMin}{10}
\newcommand{\ChThreeGeneralLastFitLag}{28440}
\newcommand{\ChThreeHangClosedFracMain}{76.5}
\newcommand{\ChThreeHangClosedFracStart}{64.0}
\newcommand{\ChThreeHangCohortOverlapPct}{99.7}
\newcommand{\ChThreeHangGeneralGroups}{7}
\newcommand{\ChThreeHangGeneralHurst}{0.841 [0.816, 0.855]}
\newcommand{\ChThreeHangGeneralReplicates}{998}
\newcommand{\ChThreeHangGeneralRms}{0.075}
\newcommand{\ChThreeHangGeneralRmsRatio}{1.19}
\newcommand{\ChThreeHangKurtEEnd}{4.03}
\newcommand{\ChThreeHangKurtEStart}{-0.60}
\newcommand{\ChThreeHangKurtNEnd}{2.24}
\newcommand{\ChThreeHangKurtNStart}{-0.72}
\newcommand{\ChThreeHangLastSupportLag}{43200}
\newcommand{\ChThreeHangLocalMinH}{0.57}
\newcommand{\ChThreeHangLocalMinLag}{320}
\newcommand{\ChThreeHangMomentContrast}{-0.0085 [-0.0112, -0.0057]}
\newcommand{\ChThreeHangMomentRmsMax}{0.126}
\newcommand{\ChThreeHangMomentRmsMin}{0.011}
\newcommand{\ChThreeHangPriorInterp}{9,648}
\newcommand{\ChThreeHangQuantRatioEnd}{0.233}
\newcommand{\ChThreeHangQuantRatioMax}{0.385}
\newcommand{\ChThreeHangQuantRatioMaxLag}{1970}
\newcommand{\ChThreeHangQuantRatioMin}{0.192}
\newcommand{\ChThreeHangQuantRatioMinLag}{90}
\newcommand{\ChThreeHangQuantRatioStart}{0.472}
\newcommand{\ChThreeHangQuantRmsMax}{0.101}
\newcommand{\ChThreeHangQuantRmsMin}{0.037}
\newcommand{\ChThreeMainClosed}{28,234}
\newcommand{\ChThreeMainClosedPct}{61.0}
\newcommand{\ChThreeMainOpen}{17,961}
\newcommand{\ChThreeMainOpenPct}{38.8}
\newcommand{\ChThreeMainUnknown}{78}
\newcommand{\ChThreeOpenAltShiftMax}{0.064}
\newcommand{\ChThreeOpenAltShiftMin}{0.051}
\newcommand{\ChThreeParaClosedFracMain}{61.0}
\newcommand{\ChThreeParaClosedFracStart}{56.6}
\newcommand{\ChThreeParaCohortOverlapPct}{99.6}
\newcommand{\ChThreeParaGeneralHurst}{0.869 [0.864, 0.874]}
\newcommand{\ChThreeParaGeneralRms}{0.038}
\newcommand{\ChThreeParaGeneralRmsRatio}{1.09}
\newcommand{\ChThreeParaKurtEEnd}{2.77}
\newcommand{\ChThreeParaKurtEStart}{-0.68}
\newcommand{\ChThreeParaKurtNEnd}{0.94}
\newcommand{\ChThreeParaKurtNStart}{-0.97}
\newcommand{\ChThreeParaLastSupportLag}{32690}
\newcommand{\ChThreeParaLocalMinH}{0.66}
\newcommand{\ChThreeParaLocalMinLag}{110}
\newcommand{\ChThreeParaMomentContrast}{0.0095 [0.0086, 0.0105]}
\newcommand{\ChThreeParaMomentRmsMax}{0.079}
\newcommand{\ChThreeParaMomentRmsMin}{0.010}
\newcommand{\ChThreeParaPriorInterp}{355,170}
\newcommand{\ChThreeParaQuantRatioEnd}{0.205}
\newcommand{\ChThreeParaQuantRatioMax}{0.398}
\newcommand{\ChThreeParaQuantRatioMaxLag}{1970}
\newcommand{\ChThreeParaQuantRatioMin}{0.244}
\newcommand{\ChThreeParaQuantRatioMinLag}{100}
\newcommand{\ChThreeParaQuantRatioStart}{0.510}
\newcommand{\ChThreeParaQuantRmsMax}{0.084}
\newcommand{\ChThreeParaQuantRmsMin}{0.023}
\newcommand{\ChThreeTaskContrastMax}{0.129}
\newcommand{\ChThreeTaskContrastMin}{0.023}
\newcommand{\ChThreeTaskGapHigh}{0.023 [0.021, 0.024]}
\newcommand{\ChThreeTaskGapHills}{0.129 [0.125, 0.132]}
\newcommand{\ChThreeTaskGapLow}{0.048 [0.046, 0.049]}
\newcommand{\ChThreeTaskGapPlains}{0.128 [0.120, 0.138]}

% Generated values only; prose is maintained in the chapter.
\newcommand{\ChThreeHangGridGroups}{267}
\newcommand{\ChThreeHangGridHFactor}{2.00--2.05}
\newcommand{\ChThreeHangGridHThreeMonths}{1.40}
\newcommand{\ChThreeHangGridMaxPct}{13.4}
\newcommand{\ChThreeHangGridMedian}{2}
\newcommand{\ChThreeHangGridMsdShiftRange}{0.98--1.87}
\newcommand{\ChThreeHangGridOldSingletonPct}{81.3}
\newcommand{\ChThreeHangGridSingletonFlightPct}{5.0}
\newcommand{\ChThreeHangGridSingletonPct}{43.4}
\newcommand{\ChThreeParaGridGroups}{1901}
\newcommand{\ChThreeParaGridHFactor}{3.97--4.05}
\newcommand{\ChThreeParaGridHThreeMonths}{1.49}
\newcommand{\ChThreeParaGridMaxPct}{4.4}
\newcommand{\ChThreeParaGridMedian}{3}
\newcommand{\ChThreeParaGridMsdShiftRange}{0.98--1.04}
\newcommand{\ChThreeParaGridOldSingletonPct}{71.1}
\newcommand{\ChThreeParaGridSingletonFlightPct}{1.4}
\newcommand{\ChThreeParaGridSingletonPct}{33.5}

% Generated by check_origin_dependence.py
\newcommand{\ChThreeHangOriginDeltaH}{-0.0168}
\newcommand{\ChThreeHangOriginDeltaHHigh}{-0.0150}
\newcommand{\ChThreeHangOriginDeltaHLow}{-0.0186}
\newcommand{\ChThreeHangOriginPeak}{1.337}
\newcommand{\ChThreeHangOriginPeakInterval}{1.310--1.364}
\newcommand{\ChThreeHangOriginPeakLag}{390}
\newcommand{\ChThreeHangOriginPeakLever}{11489}
\newcommand{\ChThreeHangOriginResolved}{32}
\newcommand{\ChThreeParaOriginDeltaH}{-0.0104}
\newcommand{\ChThreeParaOriginDeltaHHigh}{-0.0098}
\newcommand{\ChThreeParaOriginDeltaHLow}{-0.0110}
\newcommand{\ChThreeParaOriginPeak}{1.388}
\newcommand{\ChThreeParaOriginPeakInterval}{1.378--1.398}
\newcommand{\ChThreeParaOriginPeakLag}{390}
\newcommand{\ChThreeParaOriginPeakLever}{12100}
\newcommand{\ChThreeParaOriginResolved}{33}



We measure the mean-square displacement (MSD), its dependence on circuit type and
initial altitude, and the scaling of the displacement distribution. Paragliders and
hang gliders are analysed separately.

\section{Mean-square displacement}
\label{sec:fixed-msd}
\subsection{Estimator and temporal resolution}

Let $\boldsymbol{x}_f(t)=(E_f(t),N_f(t))$ be the cleaned horizontal trajectory of flight
$f$, and let $s$ index its retained continuous segments. Within each segment, define
$\Delta\boldsymbol{x}_{fi}(\tau)=\boldsymbol{x}_f(t_i+\tau)-\boldsymbol{x}_f(t_i)$
and $R_{fi}(\tau)=|\Delta\boldsymbol{x}_{fi}(\tau)|$, defined whenever $t_i$ and
$t_i+\tau$ lie in the same segment $s$; no increment crosses a segment boundary.
Let $i$ index grid origins across the flight, and define the admissible index set
\[
 \mathcal I_f(\tau)=\{i:\ t_i\text{ and }t_i+\tau
 \text{ lie in the same retained segment of flight }f\}.
\]
Let $n_{f,s}(\tau)$ count the admissible
indices in segment $s$, so that $n_f(\tau)=|\mathcal I_f(\tau)|=\sum_s n_{f,s}(\tau)$
is their total across the flight. Set
\[
 \mathcal C(\tau)=\{f : n_f(\tau)>0\}, \qquad N(\tau)=|\mathcal C(\tau)|.
\]
Let $m_f(\tau)$ be the time-averaged MSD of flight $f$, and
$M_2(\tau)$ the mean of $m_f(\tau)$ over $\mathcal C(\tau)$. They are defined by
\begin{equation}
 m_f(\tau)=\frac{1}{n_f(\tau)}\sum_{i\in\mathcal I_f(\tau)}R_{fi}(\tau)^2,
 \qquad M_2(\tau)=\frac{1}{N(\tau)}\sum_{f\in\mathcal C(\tau)}m_f(\tau).
 \label{eq:fixed-msd}
\end{equation}
We first average the squared displacements within each flight, pooling the valid
increments from all its retained segments. We then average these flight-level MSDs,
giving each contributing flight the same weight $1/N(\tau)$, regardless of its
number of increments or segments. This is the \emph{equal-flight} average.
Two alternatives were rejected.
\begin{itemize}
 \item \emph{Equal-increment average}: pooling all increments directly, without
 first averaging within flights,
 \[
  M_2^{\mathrm{incr}}(\tau)=\frac{1}{\sum_{g\in\mathcal C(\tau)}n_g(\tau)}
  \sum_{f\in\mathcal C(\tau)}\sum_{i\in\mathcal I_f(\tau)}R_{fi}(\tau)^2
  =\sum_{f\in\mathcal C(\tau)}\frac{n_f(\tau)}{\sum_{g\in\mathcal C(\tau)}n_g(\tau)}
  m_f(\tau),
 \]
 would give flight $f$ weight $n_f(\tau)/\sum_{g\in\mathcal C(\tau)}n_g(\tau)$ in
 the population TA-MSD: longer flights, which accumulate more valid increments,
 would dominate the pooled statistic, which would then track long flights rather
 than represent the population evenly.
 \item \emph{Equal-segment average}: first computing a TA-MSD for each segment,
 then averaging over every segment in $\mathcal C(\tau)$,
 \[
  m_{f,s}(\tau)=\frac{1}{n_{f,s}(\tau)}\sum_{i\in s}R_{fi}(\tau)^2, \qquad
  M_2^{\mathrm{seg}}(\tau)=\frac{1}{S(\tau)}\sum_{f\in\mathcal C(\tau)}\sum_{s}
  m_{f,s}(\tau),
 \]
 with $S(\tau)$ the total number of segments with $n_{f,s}(\tau)>0$ across
 $\mathcal C(\tau)$, would give equal weight to every segment regardless of its
 length: flights broken into many short segments, typically those with lower
 signal quality, would then be overrepresented.
\end{itemize}

For $10\leq\tau\leq10000$ s we subsample each segment to a 10 s grid, starting at its
first fix. This spacing sets the step between the origins $t_i$ from which the MSD
at lag $\tau$ is computed within a segment, and it gives every flight the same
physical lags to compare. Records with native cadence above 10 s are excluded.
Missing positions are linearly interpolated within segments: with an 8 s cadence,
for example, the 10 s position lies between the fixes at 8 and 16 s.
Table~\ref{tab:fixed-interpolation} counts how often the grid requires this. The 33
evaluated lags are multiples of 10 s, approximately logarithmically spaced and
including 100 and 1000 s. Figure markers show the evaluations; lines connect them.

The spacing itself is a trade-off. Finer sampling would raise the computational cost
sharply without benefit once it resolves timescales below any actual change in the
dynamics; coarser sampling would fail to resolve the fastest manoeuvres, which occur
during the search and thermalling phases, where speed, acceleration and heading vary
faster than during gliding. Vilpellet et al.~\cite{vilpellet2026} report that a
single revolution while thermalling takes just under 30 s, so 10 s samples each turn
about three times, resolving the fastest dynamics in the archive without
oversampling the slower gliding phase.

\begin{table}[htbp]
 \centering\small
 \begin{tabular}{lrrrr}
\toprule
Discipline & Grid fixes & Coincident & Interpolated & Fraction\\
\midrule
Paragliders & 85,915,677 & 81,888,079 & 4,027,598 & 4.69\%\\
Hang gliders & 4,099,322 & 3,639,355 & 459,967 & 11.22\%\\
\bottomrule
\end{tabular}

 \caption{Common-grid positions in the main fixed cohort. Coincidence is checked
 against actual cleaned timestamps. The last two columns count interpolation added
 by the 10 s grid, separately from reconstruction during cleaning.}
 \label{tab:fixed-interpolation}
\end{table}
Among the coincident grid positions, \ChThreeParaPriorInterp{} paraglider fixes and
\ChThreeHangPriorInterp{} hang-glider fixes already carry the cleaning interpolation
flag. The short-lag detail at 1--9 s in the general TA-MSD uses only native segments that
resolve the requested lag exactly.

\subsection{General transport and uncertainty}
\label{sec:fixed-general}

Figure~\ref{fig:fixed-general} compares the TA-MSD, averaged over within-segment
origins, with squared distance from the flight's single post-trim origin
(Section~\ref{sec:trimming}). For the latter, $t$ is elapsed time from that origin.
Gaps leave missing support without resetting either the clock or the reference
position.

A flight contributes at $t$ if a retained segment covers that time and its native
sampling interval satisfies $dt\leq t$. Linear interpolation misses any curvature
between consecutive fixes and its position error depends on motion within the local
interval $dt$, so it doesn't improve as elapsed time $t$ increases.
What can decrease is the error \emph{relative to the displacement from the origin}.
At a characteristic speed $v$, unresolved motion has a spatial scale $v\,dt$,
whereas approximately ballistic displacement grows as $v\,t$. Their ratio $dt/t$
explains the cutoff: for $t<dt$, interpolation error can be comparable to the
displacement; as $t$ grows beyond $dt$, its relative contribution can decrease.


At each $t$ we report the mean, median and 5--95\% range across available flights.
The contributing population changes as $t$ grows, since flights whose retained span
is shorter than $t$ drop out, and the TA-MSD loses flights the same way as $\tau$
grows; Table~\ref{tab:fixed-support} counts the flights that remain. These curves
therefore describe a changing mixture of sites, days and flight
conditions, which is why the controlled comparisons of
Section~\ref{sec:fixed-population} instead fix the population once for all lags.

\begin{figure}[p]
 \centering
 \includegraphics[width=\textwidth]{generated/ch3_transport_general.pdf}
 \caption{General horizontal transport. Left: squared distance from the post-trim
 origin, with equal-flight mean, median and the empirical 5--95\% flight range.
 Right: equal-flight TA-MSD over the population available at each lag, with
 pointwise 90\% cluster-bootstrap intervals and a log--log fit over
 \ChThreeGeneralFitLagMin--\ChThreeGeneralFitLagMax{} s (evaluated lags:
 \ChThreeGeneralFitLagMin--\ChThreeGeneralLastFitLag{} s). Native short-lag support
 differs from common-grid support.}
 \label{fig:fixed-general}
\end{figure}

\begin{figure}[p]
 \centering
 \includegraphics[width=\textwidth]{generated/ch3_transport_general_slopes.pdf}
 \caption{Local slopes of the general curves. The origin-distance mean and median
 have no confidence band here; the TA-MSD band is obtained by refitting bootstrap
 curves. The horizontal coordinate is elapsed time $t$ from trimming for the
 origin-distance curves and lag $\tau$ for the TA-MSD. The dotted reference is
 ballistic growth, $H=1$.}
 \label{fig:fixed-general-slopes}
\end{figure}
We summarise the growth of the TA-MSD by an analytical linear fit in logarithmic scale ($\log_{10}M_2$ against
$\log_{10}\tau$). A power law $M_2\propto\tau^{2H}$ appears in this plane as a line of
slope $2H$, so we define $H$ as half the fitted slope.\\
It's imporant to highlight that here $H$ denotes an effective growth exponent: its identification with a process's
self-similarity exponent is tested in Section~\ref{sec:fixed-scaling}.  The
interval for $H$ comes from repeating the fit on the TA-MSD recomputed for each
bootstrap replicate, a resampling of site--day groups described in the paragraph
``\nameref{par:fixed-uncertainty}'' below.
The fits to the TA-MSD of the available population (Figure~\ref{fig:fixed-general},
right) are requested over
\ChThreeGeneralFitLagMin--\ChThreeGeneralFitLagMax{} s but the last lag of the grid inside
that range is \ChThreeGeneralLastFitLag{} s, so they span
\ChThreeGeneralFitLagMin--\ChThreeGeneralLastFitLag{} s.
In the right panels of Figure~\ref{fig:fixed-general}, the TA-MSD over the fitted lags
lies close to a straight line in logarithmic scale. Figure~\ref{fig:fixed-general-residuals}
shows the fit residuals: they
fall on both sides of zero over the whole fitted range, and their RMS say that a
typical measured value lies within a factor \ChThreeParaGeneralRmsRatio{} and
\ChThreeHangGeneralRmsRatio{} of the fitted line, respectively for paragliders and hang gliders.
Such ratios, close to
unity over a curve that grows by about six decades, combined with the fact that the residuals
are evenly distributed above and below the fitted curve, make the line an adequate
description of the measured growth. The residuals are nevertheless not scattered at random.
In fact, they change sign three times as
the lag grows, so the departure from a single power law is a smooth curvature, largest
at the longest lags where the contributing population has shrunk most
(Table~\ref{tab:fixed-support}). The fit therefore summarises the growth as one
effective exponent, and the local exponent below (Figure~\ref{fig:fixed-general-slopes})
shows how the slope varies with lag.

\begin{figure}[htbp]
 \centering
 \includegraphics[width=\textwidth]{generated/ch3_transport_general_residuals.pdf}
 \caption{Residuals of the log--log fit of the general TA-MSD at the fitted lags, in dex:
 $\log_{10}M_2$ minus the fitted line. A residual of $r$ dex means the measured curve
 sits a factor $10^{r}$ from the line, so $0$ dex is a point lying on it. Dotted lines
 mark plus and minus the RMS residual.}
 \label{fig:fixed-general-residuals}
\end{figure}

The available-population fits give $H=\ChThreeParaGeneralHurst$ for paragliders and
$H=\ChThreeHangGeneralHurst$ for hang gliders and, since the flights contributing at each lag
change as $\tau$ grows (Table~\ref{tab:fixed-support}), these values summarise the
growth of a curve whose population shrinks with lag.

\begin{table}[htbp]
 \centering\small
 \begin{tabular}{lrrr}
\toprule
Discipline & Lag (s) & TA-MSD flights & Origin-distance flights\\
\midrule
Paragliders & 10 & 155,085 & 152,771\\
Paragliders & 100 & 155,044 & 155,117\\
Paragliders & 1000 & 154,482 & 155,250\\
Paragliders & 10000 & 46,273 & 72,160\\
Paragliders & 32690 & 27 & 1,237\\
Hang gliders & 10 & 6,060 & 5,330\\
Hang gliders & 100 & 6,060 & 6,059\\
Hang gliders & 1000 & 6,044 & 6,070\\
Hang gliders & 10000 & 2,326 & 3,480\\
Hang gliders & 43200 & 1 & 3\\
\bottomrule
\end{tabular}

 \caption{Number of contributing flights in the general curves. The last row for
 each discipline is the largest evaluated lag with TA-MSD support. The origin-distance
 curve uses elapsed time; the TA-MSD requires a within-segment lag no larger than
 $0.8T_s$. Their support sets therefore differ. Early origin-distance support can increase
 when elapsed time exceeds an initial excluded interval or resolves a coarse cadence.}
 \label{tab:fixed-support}
\end{table}

\Needspace{20\baselineskip}
To see how the growth changes with lag, we also compute a local exponent. For a
second-moment curve it is
\begin{equation}
 H_{\rm loc}(\tau)=\frac12\frac{d\log M_2(\tau)}{d\log\tau}.
 \label{eq:fixed-local-h}
\end{equation}
We estimate this derivative with the same closed-form least-squares slope, divided by
two, computed at each lag on the points within $\pm0.25$ decades of it, provided the
window contains at least three points.
In the general curves, a sparse interior window expands to the three nearest
log-lags: the estimate at 20 s uses the values at 10, 20 and 30 s.
The band of the TA-MSD in Figure~\ref{fig:fixed-general-slopes} comes from repeating
this local fit on every bootstrap replicate of the curve: at each lag it is the
pointwise 5--95\% range of the resulting local slopes.
For the origin-distance mean and median, we apply the same local slope to each
squared-distance curve as a function of elapsed time $t$.

Figures~\ref{fig:fixed-general} and~\ref{fig:fixed-general-slopes} show two linked
features of the origin-distance curves. After the approximately ballistic regime
at $t\lesssim10$ s, the mean rises visibly above the median, with a marked
separation at $t\sim10^2$--$10^3$ s. The mean's local slopes reach minima at \ChThreeParaLocalMinLag{} s
($H_{\rm loc}=\ChThreeParaLocalMinH$) for paragliders and \ChThreeHangLocalMinLag{} s
($H_{\rm loc}=\ChThreeHangLocalMinH$) for hang gliders within 20--2000 s, then recover
towards $H_{\rm loc}\simeq1$ around 600--700 s.


A possible physical explanation is the asynchronous alternation of gliding,
search and climb. At the same elapsed time, some pilots glide while others search
for lift or circle in a thermal, making less net horizontal progress. At short
elapsed times, an episode lasting a few hundred seconds can greatly reduce the
displacement relative to a flight that keeps gliding. The latter pulls the mean
above the median, especially because displacement is squared. Meanwhile, the
slow horizontal progress of flights in search and climb lowers the growth rates,
with a stronger dip in the median's $H_{\rm loc}$.

The duration distributions in Vilpellet et al.~\cite[Fig.~4f,i]{vilpellet2026}
support this timescale: search times and climb times have survival probabilities that remain of order $10^{-3}$ at
durations of several hundred seconds. An archive with many thousands of flights
can therefore contain several such long events alongside much shorter ones.

Search and climb continue throughout the flight, at different times for different
pilots. At $t\sim10^3$--$10^4$ s, however, many flights have already reached much
larger distances from the origin. The extra displacement gained by gliding rather
than searching or climbing during one interval of $10^2$--$10^3$ s can then be
smaller relative to that distance, and so have a smaller relative effect on
squared distance. The same alternation of phases can thus produce a weaker contrast between
flights at long elapsed times. This is consistent with the recovery of the local
exponents and the narrowing of the early mean--median gap, even though the phases
remain asynchronous.

\paragraph{Uncertainty.}\label{par:fixed-uncertainty} The shading around the origin-distance curves of
Figure~\ref{fig:fixed-general} is the empirical 5--95\% range across flights, a
measure of scatter among flights. Every other shaded band in this chapter is a
pointwise 90\% bootstrap interval. Flights are grouped by
catalogue take-off site (the one declared in the FFVL catalogue metadata
(Section~\ref{sec:catalog}, Table~\ref{tab:catalogfields})) and calendar date, retaining dependence within each site--day.
Missing keys give singleton groups.
A \emph{replicate} is one resampled version of the archive. Writing $G$ for the number of
site--day groups in the eligible archive, a replicate draws $G$ of those groups with
replacement, so a group may come out several times or not at all. We build $B=1000$
replicates. They are drawn once, over the whole eligible archive; an analysis of a
sub-population, such as a duration cohort, reuses the same draws and keeps only the
flights that belong to it.
Let $g(f)$ be the group of flight $f$, and $c^{(b)}_g$ the number of times group $g$ is
drawn in replicate $b$. The replicate replaces the equal weight $1/N(\tau)$ of
Eq.~\eqref{eq:fixed-msd} by a weight proportional to the multiplicity of the flight's
group:
\begin{equation}
 M_2^{(b)}(\tau)=\frac{1}{N^{(b)}(\tau)}
 \sum_{f\in\mathcal C(\tau)}c^{(b)}_{g(f)}\,m_f(\tau),
 \qquad
 N^{(b)}(\tau)=\sum_{f\in\mathcal C(\tau)}c^{(b)}_{g(f)}.
 \label{eq:fixed-bootstrap}
\end{equation}
Setting every $c^{(b)}_g=1$ recovers Eq.~\eqref{eq:fixed-msd}. A group drawn twice
counts its flights twice; a group never drawn contributes nothing.
The multiplicities are drawn once and then held fixed: the same $c^{(b)}_g$ serve every
lag, observable and cohort comparison, which keeps the replicates paired. Only the
normalisation $N^{(b)}(\tau)$ changes with lag, because $\mathcal C(\tau)$ shrinks as
$\tau$ grows. Quantiles, ratios and fits are recomputed in each replicate. For an estimate $\widehat\theta$,
the 5th and 95th percentiles of the bootstrap replicates $\{\widehat\theta^{(b)}\}_{b=1}^{1000}$ give a pointwise 90\% confidence interval
\begin{equation}
 I_{90}(\theta)=
 \left[\operatorname{perc}_{5}\{\widehat\theta^{(b)}\}_{b=1}^{1000},\,
       \operatorname{perc}_{95}\{\widehat\theta^{(b)}\}_{b=1}^{1000}\right].
 \label{eq:fixed-ci}
\end{equation}
Percentiles are linearly interpolated between ordered replicates.
Furthermore, at a given lag, a replicate whose resampled flights provide no support there
yields no estimate and is dropped from that lag's percentiles.

Cluster bootstrap requires many approximately independent groups and sufficiently
stable statistics~\cite{cameron_miller2015}; flights within a group need not be
independent. Table~\ref{tab:fixed-clusters} describes their spatial and temporal
separation. Nearby sites and successive days can share weather, so these distances
do not verify independence; the paragraph below checks MSD and fitted-exponent errors using geographic,
altitude and seasonal groups. Residual dependence may make the bands
too narrow; upload selection also limits inference to the population represented
by the archive. General TA-MSD bands are omitted below 20 contributing groups, an
operational support threshold rather than a guarantee of coverage.

The statistical validity of these intervals is therefore not established at this
stage. Independence between site--day groups is assumed and has never been verified,
and the size of the groups makes the assumption doubtful: their median size is one
flight, the 95th percentile is five flights for paragliders and three for hang
gliders, and even the largest group holds 47 and 13 flights
(Table~\ref{tab:fixed-clusters}). A partition this fine absorbs little of the
dependence it was introduced to capture. Flights that share one weather episode
across neighbouring sites, or across consecutive days at the same site, fall into
different groups and are resampled as if they carried independent information. The
groups are then probably not independent, and Eq.~\eqref{eq:fixed-ci} probably
understates the true uncertainty, so every band reported in this chapter should be
read as a lower bound on the spread. The open question is whether bootstrap groups
coarse enough to be approximately independent can be built at all. The paragraph
below is a first attempt at it and has not answered it yet, which is why it is
typeset in red.

\begin{table}[htbp]
 \centering\small\setlength{\tabcolsep}{4pt}
 For paragliders, the main cohort spans 3,424 catalogue sites and 3,642 dates; group size has median 1, the 95th percentile 5, and the largest group 47 (0.10\% of flights). Where another site is represented on the same day (24,509 of 25,682 groups), its nearest distance has median 11.2 km and 10--90\% range 0.04--173 km. The gap from the previous observed date at the same site is 19 days (10--90\%: 1--429; 22,258 comparisons). Missing group keys affect 0 main-cohort flights.

For hang gliders, the main cohort spans 284 catalogue sites and 1,017 dates; group size has median 1, the 95th percentile 3, and the largest group 13 (0.56\% of flights). Where another site is represented on the same day (1,147 of 1,768 groups), its nearest distance has median 24.7 km and 10--90\% range 0.04--338 km. The gap from the previous observed date at the same site is 18 days (10--90\%: 1--382; 1,484 comparisons). Missing group keys affect 0 main-cohort flights.

 \caption{Site--day diagnostics in the main cohort. Distances refer to the nearest
 other catalogue site represented on the same date; time gaps refer to the previous
 observed date at the same site. Distances and time gaps are reported as median
 [10th, 90th percentile].
 Cohort group counts are in Table~\ref{tab:fixed-cohorts}.}
 \label{tab:fixed-clusters}
\end{table}

\Needspace{12\baselineskip}
\begin{wipblock}
\paragraph{Reliability of bootstrap uncertainty estimates (work in progress, not
working).}
\label{sec:fixed-bootstrap-reliability}

This paragraph is a working draft and no conclusion in it is settled. The grouping
tried here does not yet produce approximately independent bootstrap groups, so its
intervals validate nothing and its numbers are reported only to record where the
attempt stands.

We repeat the uncertainty check with groups defined by a fixed geographic cell,
launch-altitude class and month of the year, pooling all years. Thus July flights
from different years share a group when their initial positions lie in the same
cell and altitude class. The grid is independent of the observed sites: each
initial position belongs to one half-open $50\times50$ km square in its WGS84 UTM
zone, with sides measured in projected metres. Cells are clipped at zone and
hemisphere boundaries. The four altitude classes (plains, hills, low and high
mountains) retain the existing thresholds of 300, 800 and 1500 m; they are a proxy
for launch orography. This construction groups geographical and seasonal context,
not necessarily a common weather episode.

All main-cohort flights and their equal-flight weights are retained. We draw 1000
bootstrap replicates, paired across lags, and repeat with two seeds. The archived
site--day bootstrap provides the reference. Table~\ref{tab:fixed-grid-groups}
compares group composition. Singleton groups remain, but contain only
\ChThreeParaGridSingletonFlightPct\% of paraglider flights and
\ChThreeHangGridSingletonFlightPct\% of hang-glider flights. The standard error of
$H$ is \ChThreeParaGridHFactor{} and \ChThreeHangGridHFactor{} times the site--day
value, respectively; point estimates are unchanged.

\begin{table}[htbp]
 \centering\small\setlength{\tabcolsep}{4pt}
 \begin{tabular}{lrrrrr}
\toprule
 & Site--day & Grid & Median size & Singletons (\%) & Max. share (\%) \\
\midrule
Paragliders & 25682 & 1901 & 3 & 33.5 & 4.4 \\
Hang gliders & 1768 & 267 & 2 & 43.4 & 13.4 \\
\bottomrule
\end{tabular}

 \caption{Bootstrap groups of the main cohort before and after cell--altitude--month
 grouping. ``Before'' are the site--day groups that resample every other interval in
 this chapter; ``after'' are the alternative groups introduced in this check, used only
 here. Median size is the number of flights per new group. Singletons are expressed as
 a percentage of new groups; maximum share is the largest group's percentage of
 cohort flights. All years are pooled within each month.}
 \label{tab:fixed-grid-groups}
\end{table}

We test sensitivity using 10, 25 and 50 km cells, shifting each grid by half a cell
along either or both axes. Separately, we merge two or three adjacent months and
shift the seasonal boundaries, retaining 50 km cells and the altitude classes
(Fig.~\ref{fig:fixed-bootstrap-reliability}). These are alternative resampling
partitions, not exclusions of flights.

\begin{figure}[htbp]
 \centering
 \includegraphics[width=\textwidth]{generated/ch3_transport_grid_bootstrap.pdf}
 \caption{Standard errors relative to the unshifted 50 km, one-month grid baseline.
 Left: cell size and boundary shifts. Right: adjacent-month unions with seasonal
 boundary shifts. Markers are medians and bars span minima and maxima across
 partitions and two seeds; they are not confidence intervals. Altitude classes
 remain separate, and all years are pooled.}
 \label{fig:fixed-bootstrap-reliability}
\end{figure}

The uncertainty has not stabilised across these choices. Merging three months
multiplies the median standard error of $H$ by \ChThreeParaGridHThreeMonths{} for
paragliders and \ChThreeHangGridHThreeMonths{} for hang gliders. For the latter,
shifting the 50 km grid gives MSD errors at 1000 s ranging from
\ChThreeHangGridMsdShiftRange{} times the unshifted value. Larger groups therefore
do not by themselves validate independence or interval coverage. Neighbouring
cells, months and altitude classes can still share conditions, while persistent
geographical differences can also affect the comparison~\cite{cameron_miller2015}.
The new intervals are conditional on this grouping; the archived site--day bands
elsewhere remain the reference and are not replaced by this MSD/$H$ check.
\end{wipblock}

\FloatBarrier

\subsection{Fixed population and the effect of duration selection}
\label{sec:fixed-population}

The general TA-MSD of Sec.~\ref{sec:fixed-general} averages over whichever flights
support each lag, so its population shrinks as $\tau$ grows: short flights and short
segments drop out one after another. The mix of circuit types and launch sites that
remains then changes with $\tau$, and a change in the slope of $M_2(\tau)$ cannot be
told apart from a change in the dynamics. Every statistic of the motion in the rest of
this chapter is therefore measured on one population that does not depend on $\tau$.
Such a population is limited by its shortest member: the longest lag it supports is
fixed by how long its segments are. Selecting only long segments gives more lags, and
loses flights. This subsection asks whether that selection changes the measured $H$
enough to matter, and chooses the cohort accordingly.

\paragraph{Cohorts.}
We select the segments once, using
\begin{equation}
 \tau_{\max}\leq0.8T_s,\qquad
 \mathcal C_{\tau_{\max}}=\{f:\exists s\in f,\ T_s\geq\tau_{\max}/0.8\},
 \label{eq:fixed-support}
\end{equation}
where $T_s$ is the segment's elapsed duration on the 10 s grid. The condition
$\tau_{\max}/T_s\leq0.8$ leaves $T_s-\tau_{\max}\geq0.2\,T_s$, one fifth of the
segment, for the origins at the largest lag. The threshold 0.8 fixes how many origins
remain there. For the main cohort, $\mathcal C_{10000}$, segments last at least
12500 s, so each supplies at least 251 origins on the 10 s grid at $\tau=10000$ s;
requiring only $T_s\geq\tau_{\max}$ would leave one. If a flight has segments of
14000 and 3000 s, only the first contributes, even at $\tau=10$ s. Flights and
selected segments therefore remain the same at every lag.

\paragraph{Weights.}
For each cohort, $\mathcal I_f(\tau)$ and $n_f(\tau)$ refer only to its selected
segments. Each flight of $\mathcal C_{10000}$ carries the same weight $1/N$ at every
lag, because the cohort does not change with $\tau$. The number of admissible
origins $n_f(\tau)$ still decreases with lag, so a flight's weight is divided evenly
among them. The weight of origin $i$ of flight $f$ at lag $\tau$ is
\begin{equation}
 w_{fi}(\tau)=\frac{1}{Nn_f(\tau)},\qquad
 \sum_{i\in\mathcal I_f(\tau)}w_{fi}(\tau)=\frac1N,\qquad N=|\mathcal C_{10000}|.
 \label{eq:fixed-weights}
\end{equation}
The origins of one flight therefore sum to that flight's share $1/N$, and the origins of
all flights sum to one. This is the equal-flight average of Eq.~\eqref{eq:fixed-msd},
applied to a population that does not change with $\tau$: it gives
$M_2(\tau)=\sum_f\sum_{i}w_{fi}(\tau)R_{fi}(\tau)^2$, and the same weights define the
mean $\langle A\rangle_\tau=\sum_f\sum_{i}w_{fi}(\tau)A_{fi}$ of any per-origin quantity
$A_{fi}$, used from Eq.~\eqref{eq:fixed-kurtosis} onwards.

\begin{table}[htbp]
 \centering\small
 \begin{tabular}{llrrr}
\toprule
Discipline & $\tau_{\max}$ (s) & Flights & Segments & Groups\\
\midrule
Paragliders & 100 & 155,044 & 188,312 & 82,650\\
Paragliders & 1000 & 154,482 & 167,152 & 82,390\\
Paragliders & 10000 & 46,273 & 46,367 & 25,682\\
Hang gliders & 100 & 6,060 & 7,515 & 4,264\\
Hang gliders & 1000 & 6,044 & 6,610 & 4,253\\
Hang gliders & 10000 & 2,326 & 2,327 & 1,768\\
\bottomrule
\end{tabular}

 \caption{Nested duration cohorts and their site--day groups, the resampling
 units of Sec.~\ref{sec:fixed-general}. Minimum segment
 durations are 125, 1250 and 12500 s, rounded upward when required by the 10 s grid.
 Existing cleaning and segment-admission conditions also apply. The table gives only
 counts. The full membership of each cohort, every flight and the start and stop of
 each of its segments, is saved in a manifest file with a checksum, so that later
 analyses can reuse exactly the same population.}
 \label{tab:fixed-cohorts}
\end{table}

\paragraph{Effect of the selection on $H$.}
We assess duration selection with the nested cohorts $\mathcal C_{100}$ and
$\mathcal C_{1000}$. They contain the main cohort's segments, additional shorter
segments of its flights, and additional flights. Selecting long segments also changes
who is in the sample. The declared circuit class, open-distance flights against closed
routes such as triangles and out-and-return flights (Sec.~\ref{sec:fixed-circuit}),
shifts with the cohort: the share of closed routes rises from
\ChThreeParaClosedFracStart\% in $\mathcal C_{100}$ to \ChThreeParaClosedFracMain\% in
$\mathcal C_{10000}$ for paragliders, and from \ChThreeHangClosedFracStart\% to
\ChThreeHangClosedFracMain\% for hang gliders.

We compare the fitted $H$ of two cohorts over the same lag interval, and the
comparison is paired. Every bootstrap replicate draws the same site--day groups
(Eq.~\eqref{eq:fixed-bootstrap}) for all cohorts, and the difference between two
exponents is taken replicate by replicate. Sampling noise that the cohorts share
cancels in the difference, so what remains isolates the effect of the selection
(Figure~\ref{fig:fixed-cohorts}, Table~\ref{tab:fixed-cohort-fits}). The fit covers
10--100 s for all three cohorts, and 10--1000 s for $\mathcal C_{1000}$ and
$\mathcal C_{10000}$, because the broader cohorts have no support at longer lags.

\begin{figure}[htbp]
 \centering
 \includegraphics[width=\textwidth]{generated/ch3_transport_cohorts.pdf}
 \caption{Effect of duration selection. Top: MSD in the three nested cohorts;
 dashed lines fit the main cohort over 10--10000 s. Bottom: effective $H$ with 90\%
 intervals, fitted over 10--100 s for all three cohorts and 10--1000 s for the two
 longer cohorts. Fits within each comparison use identical lag values.}
 \label{fig:fixed-cohorts}
\end{figure}
\begin{table}[htbp]
 \centering\small
 \begin{tabular}{lllr}
\toprule
Discipline & Fit range (s) & Contrast & $\Delta H$ [5--95\%]\\
\midrule
Paragliders & 10--100 & $H_{\mathcal C_{100}}-H_{\mathcal C_{1000}}$ & -0.00006 [-0.00009, -0.00004]\\
Paragliders & 10--100 & $H_{\mathcal C_{100}}-H_{\mathcal C_{10000}}$ & -0.01687 [-0.01712, -0.01660]\\
Paragliders & 10--100 & $H_{\mathcal C_{1000}}-H_{\mathcal C_{10000}}$ & -0.01680 [-0.01705, -0.01654]\\
Paragliders & 10--1000 & $H_{\mathcal C_{1000}}-H_{\mathcal C_{10000}}$ & -0.01703 [-0.01733, -0.01674]\\
Hang gliders & 10--100 & $H_{\mathcal C_{100}}-H_{\mathcal C_{1000}}$ & -0.00002 [-0.00013, 0.00011]\\
Hang gliders & 10--100 & $H_{\mathcal C_{100}}-H_{\mathcal C_{10000}}$ & -0.01646 [-0.01748, -0.01552]\\
Hang gliders & 10--100 & $H_{\mathcal C_{1000}}-H_{\mathcal C_{10000}}$ & -0.01644 [-0.01746, -0.01550]\\
Hang gliders & 10--1000 & $H_{\mathcal C_{1000}}-H_{\mathcal C_{10000}}$ & -0.01771 [-0.01879, -0.01668]\\
\bottomrule
\end{tabular}

 \caption{Paired differences of the cohort exponents in Fig.~\ref{fig:fixed-cohorts}.
 Each comparison uses identical lags and bootstrap draws. A 90\% interval excluding
 zero resolves a difference under the adopted resampling; inclusion of zero alone
 would not establish absence of selection bias.}
 \label{tab:fixed-cohort-fits}
\end{table}
\paragraph{Result.}
On 10--100 s, $\mathcal C_{100}$ and $\mathcal C_{1000}$ differ in $H$ by less than
$10^{-4}$ in both disciplines. The paired interval resolves this tiny difference for
paragliders; for hang gliders it includes zero. The two cohorts share
\ChThreeParaCohortOverlapPct\% (paragliders) and \ChThreeHangCohortOverlapPct\% (hang
gliders) of their flights, so in each replicate their exponents move together and
their difference has a very narrow interval. That is why a difference this small is
resolved at all.

The main cohort behaves differently. Its effective exponent is larger than that of
every broader cohort, in both fit intervals and both disciplines, and every one of
these paired differences excludes zero. The difference between the main cohort and a
broader one lies between \ChThreeCohortShiftMin{} and \ChThreeCohortShiftMax{} in $H$,
which is \ChThreeCohortBiasMinPct--\ChThreeCohortBiasMaxPct\% of the exponent itself.

We accept this bias and adopt $\mathcal C_{10000}$ for every analysis below that
requires a fixed population. A shift of about 2\% in $H$ is small beside what the cohort
provides: a population that stays the same over three decades of lag, 10--10000 s,
against two for $\mathcal C_{1000}$ and one for $\mathcal C_{100}$. Two limits apply. The shift can be measured only
where the broader cohorts have support, on 10--100 s and 10--1000 s. And every exponent reported below is
conditional on flights with an uninterrupted segment of at least 12500 s.

\FloatBarrier

\section{Transport by circuit type and initial altitude}
\label{sec:fixed-circuit}

The catalogue's declared circuit class separates open distance flights from closed
routes such as triangles and out-and-return flights. We cross this distinction with
initial cleaned GNSS altitude: Plains ($z_0<300$ m), Hills ($300\leq z_0<800$ m),
Low mountains ($800\leq z_0<1500$ m), and High mountains ($z_0\geq1500$ m).
These are launch-altitude classes; they do not classify the terrain crossed later.
This comparison uses paragliders and the same $\mathcal C_{10000}$ support rule in
all eight cells. Figure~\ref{fig:fixed-task} shows the eight MSD curves; counts and
fitted exponents are given in Table~\ref{tab:fixed-task}.
Eligible flights: \ChThreeEligibleOpen{} open (\ChThreeEligibleOpenPct\%),
\ChThreeEligibleClosed{} closed (\ChThreeEligibleClosedPct\%), and
\ChThreeEligibleUnknown{} unclassified.

Main-cohort flights: \ChThreeMainOpen{} open (\ChThreeMainOpenPct\%),
\ChThreeMainClosed{} closed (\ChThreeMainClosedPct\%), and \ChThreeMainUnknown{}
unclassified.

\begin{figure}[htbp]
 \centering
 \includegraphics[width=\textwidth]{generated/ch3_transport_task_altitude.pdf}
 \caption{Paraglider transport by declared circuit and initial altitude, with fixed
 flights and segments throughout 10--10000 s. Markers give MSD estimates, shaded
 bands their pointwise 90\% intervals, and dashed lines the global log--log fits.
 The identical axes permit comparison of displacement amplitude and growth rate.}
 \label{fig:fixed-task}
\end{figure}
\begin{table}[htbp]
 \centering\footnotesize\setlength{\tabcolsep}{3.5pt}
 \begin{tabular}{llrrrr}
\toprule
Circuit & Initial altitude & $N$ & $G$ & $H$ [5--95\%] & $M_2(10^4)$ (km$^2$)\\
\midrule
Open & Plains & 1,679 & 1,140 & 0.953 [0.952, 0.955] & 5900\\
Open & Hills & 2,068 & 1,634 & 0.944 [0.943, 0.946] & 4920\\
Open & Low mountains & 7,412 & 5,255 & 0.893 [0.892, 0.895] & 1814\\
Open & High mountains & 6,802 & 4,794 & 0.889 [0.888, 0.890] & 1710\\
Closed & Plains & 165 & 157 & 0.825 [0.816, 0.834] & 372\\
Closed & Hills & 839 & 756 & 0.816 [0.812, 0.819] & 299\\
Closed & Low mountains & 14,649 & 7,954 & 0.846 [0.845, 0.847] & 596\\
Closed & High mountains & 12,581 & 7,406 & 0.866 [0.865, 0.867] & 1001\\
\bottomrule
\end{tabular}

 \caption{Main-cohort flights $N$, site--day groups $G$, global effective exponents
 and displacement amplitude at 10000 s. Intervals come from the coupled cluster
 bootstrap. Even the smallest cell retains 165 flights from 157 groups.}
 \label{tab:fixed-task}
\end{table}

For open routes, Plains and Hills have larger $H$ than either mountain class. Thus
the distinction involves the rate of growth as well as the displacement reached at
10000 s. Closed routes have smaller $H$ than open routes in every altitude class.
The paired contrasts $H_{\rm open}-H_{\rm closed}$ range from \ChThreeTaskContrastMin{}
to \ChThreeTaskContrastMax; all four 90\% intervals exclude zero. The two lowland
classes carry the largest contrast, close to \ChThreeTaskContrastMax, and it falls to
\ChThreeTaskContrastMin{} in High mountains.
The closed-route exponents also differ between altitude classes: the present intervals
do not support a common value across them. Their ordering reverses the open-route one.
Paired lowland-minus-mountain contrasts, each lowland class against each mountain class
within one circuit, are positive for open routes, between \ChThreeOpenAltShiftMin{} and
\ChThreeOpenAltShiftMax{} in $H$, and negative for closed routes, with magnitudes between
\ChThreeClosedAltShiftMin{} and \ChThreeClosedAltShiftMax; all eight intervals exclude
zero. Launch altitude therefore enters $H$ with a sign that depends on the declared
circuit.

One reading connects the two orderings through a single constraint, how closely the
terrain lets a declared route be flown as planned. A closed route commits the pilot to a
return: once about half the available time has elapsed, the remainder of the flight
carries the glider back towards its start, and displacements measured over lags
approaching the circuit's duration lose the advance accumulated earlier. The suppression
is strongest where the return can be flown most directly. Lowland terrain imposes no
relief to circumvent and no ridge to follow, so a lowland circuit closes more faithfully,
and its net horizontal transport is the more strongly suppressed. Figure~\ref{fig:fixed-task}
shows where this acts: over the last decade of lag every closed curve bends below its own
fitted line, while the two lowland open curves stay on theirs out to 10000 s. The same freedom
acts with the opposite sign on an open route, which asks only for distance: with no
relief to follow, the track can stay straighter, and lowland open
flights hold a nearly ballistic growth out to 10000 s while the mountain classes lose
slope. The orographic setting thus reaches $H$ through the flight plan, and the sign of
its effect is set by what the plan asks for.

Three limits attach to this reading. The declared class does not guarantee exact closure
of the trimmed trajectory, and the altitude classes describe the launch site rather than
the terrain crossed later. The two lowland closed cells are the smallest in
Table~\ref{tab:fixed-task} and have the widest intervals, although their contrasts
against the mountain cells still exclude zero. And the comparison is observational: it
does not isolate orography from weather, route selection, pilots or equipment.

Within the fixed cohort the group fractions $p_g=N_g/N$ are constant. For an
exhaustive partition, including unclassified flights,
\begin{equation}
 M_2(\tau)=\sum_g p_g M_{2,g}(\tau),\qquad
 H_{\rm loc}(\tau)=\sum_g
 \frac{p_gM_{2,g}(\tau)}{M_2(\tau)}H_{{\rm loc},g}(\tau).
 \label{eq:fixed-mixture}
\end{equation}
The derivative slope weights each group by its contribution to the second moment.
Equation~\eqref{eq:fixed-mixture} describes the main-cohort MSD; the plotted local
regressions approximate its derivatives. Both circuit classes contribute substantially
to this aggregate, with distinct growth rates.

Wind can contribute to both amplitude and apparent scaling. For an additive constant
velocity $\boldsymbol u$, write $\Delta\boldsymbol x=\Delta\boldsymbol y+\boldsymbol u\tau$:
\begin{equation}
 \langle|\Delta\boldsymbol x|^2\rangle
 =\langle|\Delta\boldsymbol y|^2\rangle
 +2\tau\boldsymbol u\cdot\langle\Delta\boldsymbol y\rangle+|\boldsymbol u|^2\tau^2.
 \label{eq:fixed-wind}
\end{equation}
Even a uniform constant drift can change the slope of this uncentred MSD through its
ballistic term. A common deterministic shift cancels from the centred covariance;
drifts differing between flights need not cancel from the mixture. A variable wind
field can additionally modify velocity correlations. Wind is not separated quantitatively
in this comparison.
\FloatBarrier

\section{Anisotropy and principal directions}
\label{sec:fixed-anisotropy}
\emph{Reserved for the analysis of anisotropy and principal directions. The existing
PCA results remain in the experiments volume; they are not updated in this chapter.}

\section{Scaling of displacement distributions}
\label{sec:fixed-scaling}

The MSD measures one moment of the displacement law. To test whether the whole law
rescales with one exponent, we use quantiles, a spectrum of moments and excess kurtosis.
All use the main cohort's flights, segments and weights in Eq.~\eqref{eq:fixed-weights}.

\subsection{Self-similarity and its observable consequences}
\label{sec:fixed-selfsimilar-definition}

A process $\boldsymbol X(t)$, with $\boldsymbol X(0)=0$, is $H$-self-similar if,
for some $H>0$,
\begin{equation}
 \bigl(\boldsymbol X(at_1),\ldots,\boldsymbol X(at_m)\bigr)
 \overset d= a^H\bigl(\boldsymbol X(t_1),\ldots,\boldsymbol X(t_m)\bigr)
 \quad(a>0).
 \label{eq:fixed-selfsimilar}
\end{equation}
This identity must hold for every $m\geq1$ and $t_1,\ldots,t_m\geq0$.
Equality in distribution concerns positions at any collection of times, including
their dependence~\cite{embrechts_maejima_selfsimilar}. Gaussianity and independence
of increments are not required.

Two separate hypotheses stand between this definition and a measurement, and each
does different work. Stationary increments make a pooled lag statistic informative
about $\boldsymbol X$ itself: with $\boldsymbol X(0)=0$ they give
$\boldsymbol X(\tau)\overset d=\boldsymbol X(t+\tau)-\boldsymbol X(t)$, so
displacements collected from many origins estimate the law of the position reached
after an elapsed time $\tau$. The archive supplies no repeated releases from a common
initial condition, which leaves this the only route from the increments we observe
to the process we describe. Self-similarity then compresses the resulting family of
laws: one exponent and one shape determine every lag, including elapsed times the
cohort does not cover, and that exponent is the quantity stochastic transport
models predict. Without it this chapter reports a separate
propagator at each measured lag, with no relation between them. Testing
Eq.~\eqref{eq:fixed-selfsimilar} is therefore a test of whether one anomalous
exponent describes these flights at all, and the sections below report necessary
consequences of that identity rather than a proof of it.

To infer a lag-only scaling law, we additionally assume stationary increments:
their joint laws must be unchanged by shifting the origin.
Self-similarity alone does not impose this condition. With stationary increments,
each amplitude $Y\in\{|\Delta E|,|\Delta N|,R\}$ obeys
\begin{equation}
 Y(\tau)\overset d=(\tau/\tau_0)^H Y(\tau_0),\qquad
 \begin{aligned}
 Q_{Y,p}(\tau)&=(\tau/\tau_0)^H Q_{Y,p}(\tau_0),\\
 M_{q,Y}(\tau)&=(\tau/\tau_0)^{qH}M_{q,Y}(\tau_0).
 \end{aligned}
 \label{eq:fixed-scale-family}
\end{equation}
Here $Q_{Y,p}$ is a population quantile and $M_{q,Y}$ a finite moment of order $q$;
below, the same symbols denote their equal-flight estimates.
Thus every quantile has the same exponent, quantile ratios are constant, and the
moment spectrum is $\zeta_Y(q)=qH$. Signed-component excess kurtosis is also constant
when the fourth moment exists and variance is nonzero; it need not vanish. These are
necessary consequences for the lag distributions. A finite selection of such
statistics cannot establish the joint-law identity in Eq.~\eqref{eq:fixed-selfsimilar}.

\subsection{A common empirical distribution}

With the weights in Eq.~\eqref{eq:fixed-weights}, the common empirical distribution is
\begin{equation}
 \widehat F_{Y,\tau}(y)=\frac1N\sum_{f=1}^N\frac1{n_f(\tau)}
 \sum_{i\in\mathcal I_f(\tau)}\mathbf1\{Y_{fi}(\tau)\leq y\},\qquad
 Q_{Y,p}(\tau)=\inf\{y:\widehat F_{Y,\tau}(y)\geq p\}.
 \label{eq:fixed-cdf}
\end{equation}
We average the within-flight fractions below $y$, then invert the resulting CDF.
The quantile is the first observed amplitude at which the cumulative weight reaches
$p$, without interpolation between ranks. This order of operations defines the
quantile of the equal-flight mixture.

Absolute components measure amplitudes along east and north. For example,
$\Pr(|\Delta E|\leq a)=\Pr(-a\leq\Delta E\leq a)$, so small absolute quantiles
measure proximity to zero displacement. The radius $R=(\Delta E^2+\Delta N^2)^{1/2}$
is invariant under rotation and measures horizontal displacement without assuming
isotropy.

\subsection{Quantiles, growth exponents and relative spread}
\label{sec:fixed-quantiles}

We evaluate $p=0.25,0.50,0.75,0.90$ and fit each quantile over the same 33 lags.
Writing $\mathcal B$ for the least-squares slope against $\log_{10}\tau$
used in Section~\ref{sec:fixed-general},
\begin{equation}
 H_{Y,p}=\mathcal B\!\left[\log_{10}\frac{Q_{Y,p}}{1\,\mathrm m}\right],\qquad
 \mathcal B\!\left[\log_{10}\frac{Q_{Y,0.25}}{Q_{Y,0.90}}\right]
 =H_{Y,0.25}-H_{Y,0.90}.
 \label{eq:fixed-ratio}
\end{equation}
No factor of two enters a length exponent. The ratio compares small and large
amplitudes at each lag, while the fitted exponents summarise their growth across
10--10000 s. The equality follows from linearity of least squares on identical lags;
ratio drift and the exponent difference are therefore linked diagnostics.
Figures~\ref{fig:fixed-quantiles-para} and~\ref{fig:fixed-quantiles-hang} give the
quantile curves and their ratio for each discipline, and
Figure~\ref{fig:fixed-quantile-h} collects the fitted quantile exponents.

\begin{figure}[H]
 \centering
 \includegraphics[width=\textwidth]{generated/ch3_transport_quantiles_para.pdf}
 \caption{Paraglider equal-flight displacement quantiles and their relative spread,
 on the fixed main cohort. Every panel uses the empirical measure in
 Eq.~\eqref{eq:fixed-cdf}. Shading gives pointwise 90\% cluster-bootstrap intervals.}
 \label{fig:fixed-quantiles-para}
\end{figure}
\begin{figure}[htbp]
 \centering
 \includegraphics[width=\textwidth]{generated/ch3_transport_quantiles_hang.pdf}
 \caption{The same amplitude quantiles and ratio for hang gliders, with an independent
 main cohort and the same sampling convention.}
 \label{fig:fixed-quantiles-hang}
\end{figure}
\begin{figure}[htbp]
 \centering
 \includegraphics[width=\textwidth]{generated/ch3_transport_quantile_exponents.pdf}
 \caption{Global quantile exponents from 10--10000 s fits. Error bars are the 5th--95th
 percentiles of the refitted bootstrap slopes. A common scale exponent would give
 a horizontal sequence within each panel.}
 \label{fig:fixed-quantile-h}
\end{figure}
For paragliders, the radial ratio falls from \ChThreeParaQuantRatioStart{} at 10 s to
\ChThreeParaQuantRatioMin{} at \ChThreeParaQuantRatioMinLag{} s, rises to
\ChThreeParaQuantRatioMax{} at \ChThreeParaQuantRatioMaxLag{} s, then falls to
\ChThreeParaQuantRatioEnd{} at 10000 s. Quantile-fit RMS residuals span
\ChThreeParaQuantRmsMin--\ChThreeParaQuantRmsMax{} dex.

For hang gliders, the radial ratio falls from \ChThreeHangQuantRatioStart{} at 10 s to
\ChThreeHangQuantRatioMin{} at \ChThreeHangQuantRatioMinLag{} s, rises to
\ChThreeHangQuantRatioMax{} at \ChThreeHangQuantRatioMaxLag{} s, then falls to
\ChThreeHangQuantRatioEnd{} at 10000 s. Quantile-fit RMS residuals span
\ChThreeHangQuantRmsMin--\ChThreeHangQuantRmsMax{} dex.

A positive global slope of these nonmonotonic ratios can coexist with a net endpoint decrease.
\FloatBarrier

\subsection{Moment spectrum and finite-range scaling}
\label{sec:fixed-moments}

For each amplitude and order we calculate
\begin{equation}
 M_{q,Y}(\tau)=\frac1N\sum_f\frac1{n_f(\tau)}\sum_{i\in\mathcal I_f(\tau)}Y_{fi}(\tau)^q,
 \qquad \zeta_Y(q)=\mathcal B\!\left[\log_{10}\frac{M_{q,Y}}{(1\,\mathrm m)^q}\right].
 \label{eq:fixed-moments}
\end{equation}
We use $q\in\{0.25,0.5,0.75,1,1.5,2,2.5,3,3.5,4\}$. Each point of the spectrum is
obtained by first computing a moment at all lags, then fitting its log--log slope.
The reference $\zeta=qH$ is fitted through the origin with equal weight for the ten
orders. Higher orders emphasise large displacements. The radial spectrum contains
the MSD, with the identity
\begin{equation}
 M_{2,R}=M_{2,|\Delta E|}+M_{2,|\Delta N|}=M_2
 \label{eq:fixed-second-identity}
\end{equation}
verified at every lag and in every bootstrap replicate.
Table~\ref{tab:fixed-fits} reports the main-cohort MSD exponents,
Figure~\ref{fig:fixed-local} compares their local and global values, and
Figure~\ref{fig:fixed-spectrum} gives the moment spectra.
Table~\ref{tab:fixed-scaling} collects the paired quantile and moment contrasts that
test $\zeta=qH$.

\begin{table}[htbp]
 \centering\small
 \begin{tabular}{llrr}
\toprule
Discipline & Fit range (s) & $H$ [5--95\%] & RMS (dex)\\
\midrule
Paragliders & 10--10000 & 0.882 [0.881, 0.883] & 0.046\\
Paragliders & 10--100 & 0.925 [0.924, 0.925] & 0.007\\
Paragliders & 100--1000 & 0.898 [0.897, 0.898] & 0.006\\
Paragliders & 1000--10000 & 0.804 [0.802, 0.807] & 0.026\\
Hang gliders & 10--10000 & 0.852 [0.849, 0.855] & 0.072\\
Hang gliders & 10--100 & 0.915 [0.914, 0.916] & 0.006\\
Hang gliders & 100--1000 & 0.873 [0.871, 0.874] & 0.007\\
Hang gliders & 1000--10000 & 0.728 [0.720, 0.737] & 0.051\\
\bottomrule
\end{tabular}

 \caption{Main-cohort MSD exponents: a single global fit and three common decade
 intervals. RMS residuals are measured in $\log_{10}M_2$. Each fit retains the same
 flight and segment population.}
 \label{tab:fixed-fits}
\end{table}
\begin{figure}[htbp]
 \centering
 \includegraphics[width=\textwidth]{generated/ch3_transport_fixed_slopes.pdf}
 \caption{Local and global effective MSD exponents on the fixed main cohort.
 Bands are pointwise 90\% bootstrap intervals. The global fit compresses the full
 curve into one number; the local slopes expose its scale dependence.}
 \label{fig:fixed-local}
\end{figure}
\begin{figure}[htbp]
 \centering
 \includegraphics[width=\textwidth]{generated/ch3_transport_spectrum.pdf}
 \caption{Moment spectra for the three displacement amplitudes. Each point is the
 slope of a 10--10000 s log--log moment fit; shading gives its pointwise 90\% interval.
 Dashed lines are constrained linear fits $\zeta=qH$.}
 \label{fig:fixed-spectrum}
\end{figure}
\begin{table}[htbp]
 \centering\small\setlength{\tabcolsep}{4pt}
 \begin{tabular}{llrr}
\toprule
Discipline & Amplitude & $H_{25}-H_{90}$ & $\zeta(4)/4-\zeta(0.25)/0.25$\\
\midrule
Paragliders & $|\Delta E|$ & 0.027 [0.025, 0.029] & 0.0113 [0.0097, 0.0128]\\
Paragliders & $|\Delta N|$ & 0.030 [0.028, 0.032] & -0.0004 [-0.0016, 0.0009]\\
Paragliders & $R$ & 0.023 [0.022, 0.024] & 0.0095 [0.0086, 0.0105]\\
Hang gliders & $|\Delta E|$ & 0.038 [0.032, 0.044] & 0.0095 [0.0040, 0.0150]\\
Hang gliders & $|\Delta N|$ & 0.101 [0.096, 0.106] & -0.0269 [-0.0305, -0.0232]\\
Hang gliders & $R$ & 0.075 [0.071, 0.079] & -0.0085 [-0.0112, -0.0057]\\
\bottomrule
\end{tabular}

 \caption{Paired global-slope contrasts, with 5--95\% bootstrap intervals. A common
 scale exponent predicts zero for both the quantile and moment contrasts.}
 \label{tab:fixed-scaling}
\end{table}
For paragliders, the radial contrast $\zeta(4)/4-\zeta(0.25)/0.25$ is
\ChThreeParaMomentContrast. Moment-fit RMS residuals span
\ChThreeParaMomentRmsMin--\ChThreeParaMomentRmsMax{} dex across all amplitudes and orders.

For hang gliders, the radial contrast $\zeta(4)/4-\zeta(0.25)/0.25$ is
\ChThreeHangMomentContrast. Moment-fit RMS residuals span
\ChThreeHangMomentRmsMin--\ChThreeHangMomentRmsMax{} dex across all amplitudes and orders.
\FloatBarrier

\subsection{Excess kurtosis of the signed components}
\label{sec:fixed-kurtosis}

Departures from Gaussian shape are assessed on the signed components: taking an
absolute value or a radius changes even a Gaussian law. With the equal-flight mean
$\langle A\rangle_\tau=\sum_f\sum_{i\in\mathcal I_f(\tau)}w_{fi}(\tau)A_{fi}$, define
\begin{equation}
 \mu_j(\tau)=\langle\Delta x_j\rangle_\tau,\qquad
 K_j(\tau)=\frac{\langle(\Delta x_j-\mu_j)^4\rangle_\tau}
 {\langle(\Delta x_j-\mu_j)^2\rangle_\tau^2}-3,\qquad j\in\{E,N\}.
 \label{eq:fixed-kurtosis}
\end{equation}
Centring uses the mean of the full mixture at each lag. A Gaussian marginal has
$K_j=0$~\cite{nist_kurtosis}; nonzero excess measures a departure in its standardised
fourth moment. Zero excess alone does not establish Gaussianity. For self-similarity,
the relevant prediction is constancy with lag, regardless of the value.
Figure~\ref{fig:fixed-kurtosis} shows $K_j(\tau)$ for both components.

\begin{figure}[htbp]
 \centering
 \includegraphics[width=\textwidth]{generated/ch3_transport_kurtosis.pdf}
 \caption{Signed-component Fisher excess kurtosis at the same lags, with the same
 flights, segments and weights as the quantiles and moments. The dashed line is the
 Gaussian marginal reference. Bands use the coupled cluster bootstrap.}
 \label{fig:fixed-kurtosis}
\end{figure}
For paragliders, $(K_E,K_N)$ changes from
$(\ChThreeParaKurtEStart,\ChThreeParaKurtNStart)$ at 10 s to
$(\ChThreeParaKurtEEnd,\ChThreeParaKurtNEnd)$ at 10000 s.

For hang gliders, $(K_E,K_N)$ changes from
$(\ChThreeHangKurtEStart,\ChThreeHangKurtNStart)$ at 10 s to
$(\ChThreeHangKurtEEnd,\ChThreeHangKurtNEnd)$ at 10000 s.
Different flight means and variances can themselves produce nonzero excess in the
mixture. The measured kurtosis therefore characterises the pooled displacement law.

\subsection{Assessment of self-similarity}

The quantile curves and the nearly straight moment spectra can be reconciled through
\begin{equation}
 M_{q,Y}(\tau)=\int_0^1 Q_{Y,p}(\tau)^q\,dp.
 \label{eq:fixed-moment-quantile}
\end{equation}
Each moment averages over probability levels; the order $q$ controls the emphasis
on larger amplitudes. Fitting its slope
across three decades further averages over crossovers. Changes in different parts
of the distribution can consequently leave the fitted spectrum nearly linear.
Moreover, Table~\ref{tab:fixed-scaling} resolves small departures from $\zeta=qH$;
the spectrum is not exactly linear under the adopted bootstrap.

For these data, the lag-dependent quantile ratios provide the more direct evidence
against a fixed distributional shape. Their large, nonmonotonic changes are visible
without fitting a power law, and their estimation does not require finite high-order
moments. Unequal fitted quantile exponents support this reading, but the curved
quantile trajectories limit their interpretation as constant exponents. The change
in kurtosis supplies complementary evidence. Visual linearity of ten fitted moment
slopes is therefore insufficient to outweigh the observed shape changes.

Over 10--10000 s, a single self-similar model with stationary increments does not
adequately describe the pooled displacement law.
Approximate scaling over narrower intervals remains possible. Fixed flights and
segments remove changing membership, and the contrast below measures how much of the
remaining structure follows from where within a segment the origins fall. The present
analysis does not determine multi-time laws. It therefore does not exclude every
nonstationary self-similar process, nor establish an asymptotic multifractal law.

% --- merged from sections/04-origin-dependence.tex ---
\Needspace{12\baselineskip}
\paragraph{Does the increment law depend on the origin?}
\label{sec:fixed-origin-dependence}

Equation~\eqref{eq:fixed-scale-family} reaches the lag distributions through
stationary increments, and the estimators above pool every admissible origin inside
a selected segment. That pooling estimates one population law only if the law of
$\boldsymbol X(t+\tau)-\boldsymbol X(t)$ does not depend on $t$. We test this
directly. Within each selected segment the admissible origins at lag $\tau$ are
split into three contiguous blocks of equal size, and the same equal-flight
amplitude statistics are measured separately in each. The early and late blocks
receive identical origin counts in every segment, so each cohort flight contributes
to both and the contrast is paired inside the site--day draws of
Sec.~\ref{sec:fixed-bootstrap-reliability}. Under stationary increments the
late-to-early ratio equals one at every lag. Pooling the three blocks returns the
equal-flight moments of Sec.~\ref{sec:fixed-scaling} exactly, which fixes the
diagnostic to the published measurement rather than to a separate estimator.

\begin{figure}[htbp]
 \centering
 \includegraphics[width=\textwidth]{generated/ch3_transport_origin_dependence.pdf}
 \caption{Origin dependence of the increment law in the main cohort. Top: the ratio
 of the late-block to the early-block radial moments, with 95\% paired bootstrap
 bands and the stationary-increment value as a dashed line. Bottom left: the mean
 separation between the two blocks of origins, which shrinks as the lag consumes the
 segment. Bottom right: $H$ fitted in each block over 10--10000 s. Bands and bars are
 pointwise.}
 \label{fig:fixed-origin-dependence}
\end{figure}

For paragliders, the late-origin second moment differs from the early-origin one at
\ChThreeParaOriginResolved{} of 33 lags. The largest departure is a ratio of
\ChThreeParaOriginPeak{} (\ChThreeParaOriginPeakInterval) at
$\tau=\SI{\ChThreeParaOriginPeakLag}{\second}$, where the two blocks are separated by
\SI{\ChThreeParaOriginPeakLever}{\second} on average. Fitting $H$ over 10--10000 s in
each block gives a late-minus-early difference of $\ChThreeParaOriginDeltaH$, with a
95\% interval from $\ChThreeParaOriginDeltaHLow$ to $\ChThreeParaOriginDeltaHHigh$.

For hang gliders, the late-origin second moment differs from the early-origin one at
\ChThreeHangOriginResolved{} of 33 lags. The largest departure is a ratio of
\ChThreeHangOriginPeak{} (\ChThreeHangOriginPeakInterval) at
$\tau=\SI{\ChThreeHangOriginPeakLag}{\second}$, where the two blocks are separated by
\SI{\ChThreeHangOriginPeakLever}{\second} on average. Fitting $H$ over 10--10000 s in
each block gives a late-minus-early difference of $\ChThreeHangOriginDeltaH$, with a
95\% interval from $\ChThreeHangOriginDeltaHLow$ to $\ChThreeHangOriginDeltaHHigh$.

The ratio departs from one over the whole lag range in both disciplines, so the
increments of these segments are not stationary. Displacements measured from late
origins exceed those from early origins across most of the range, with the largest
excess near $\tau\approx400$ s, and the excess closes at the longest lags. The two
blocks are separated by about 12000 s at small lags and by roughly half that at
$\tau=10000$ s, because the admissible origins shrink as the lag consumes the
segment; the contrast at the longest lags therefore rests on a shorter lever arm and
is not comparable in strength to the rest of the curve.

Two consequences follow for this chapter. First, the amplitude of the displacement
law is a property of the flight phase from which origins are drawn, so the pooled
one-lag laws of Sec.~\ref{sec:fixed-scaling} describe a mixture over flight phases
rather than one stationary increment law. Second, the fitted exponent is far less
sensitive to this mixture than the amplitude: the block-to-block change in $H$ is
resolved but small against the amplitude ratios in the same figure. Reading $H$ as a
lag exponent of the pooled mixture survives the test; reading it as the
self-similarity exponent of a stationary-increment position process does not.

The sign of the effect does not match waiting-time ageing, which suppresses late
displacements in a subordinated walk. Launch and initial circling confine the early
part of a flight, while route-following occupies its later part, and this ordering
alone produces the observed direction. The diagnostic separates origin position from
lag; it does not identify the mechanism, and the flight-phase, wind and pilot
explanations are not distinguished here.
\FloatBarrier
